\chapter{Relativistic Scalar Fields and Canonical Quantization}
\label{ch:relativistic-scalar-fields-canonical-quantization}
\ConventionsLorentzian


\section{Fields and Variational Equations}

Let \(\mathbb M^D\) be \(D\)-dimensional Minkowski spacetime with metric
\[
  \eta_{\mu\nu}=\operatorname{diag}(-1,+1,\ldots,+1).
\]
Write
\[
  x^\mu=(x^0,\vec x),\qquad x^0=t,\qquad \vec x\in\mathbb R^{D-1}.
\]
A real classical scalar field is a real-valued function
\[
  \phi:\mathbb M^D\to\mathbb R.
\]
The field-theoretic analogue of finitely many coordinates \(q^i(t)\) is the
family of coordinates \(\phi(t,\vec x)\), labelled continuously by
\(\vec x\). A Lagrangian field theory is specified by an action functional
\[
  S[\phi]
  =
  \int_{\mathbb M^D}\dd^D x\,
  \mathcal L\bigl(\phi(x),\partial_\mu\phi(x)\bigr),
\]
where \(\mathcal L\) is a Lorentz scalar Lagrangian density in the examples
below.

For a compactly supported variation \(\delta\phi\),
\[
\begin{aligned}
  \delta S
  &=
  \int\dd^Dx\,
  \left[
    \frac{\partial\mathcal L}{\partial\phi}\,\delta\phi
    +
    \frac{\partial\mathcal L}
         {\partial(\partial_\mu\phi)}
    \,\partial_\mu\delta\phi
  \right]  \\
  &=
  \int\dd^Dx\,\delta\phi
  \left[
    \frac{\partial\mathcal L}{\partial\phi}
    -
    \partial_\mu
    \frac{\partial\mathcal L}
         {\partial(\partial_\mu\phi)}
  \right].
\end{aligned}
\]
The Euler--Lagrange equation is therefore
\[
  \frac{\partial\mathcal L}{\partial\phi}
  -
  \partial_\mu
  \frac{\partial\mathcal L}{\partial(\partial_\mu\phi)}
  =
  0.
\]

\section{The Free Massive Scalar}

The free massive real scalar field of mass \(m\ge0\) has Lagrangian density
\[
  \mathcal L
  =
  \frac12(\partial_t\phi)^2
  -\frac12\sum_{i=1}^{D-1}(\partial_i\phi)^2
  -\frac12m^2\phi^2.
\]
Equivalently,
\[
  \mathcal L
  =
  -\frac12\partial_\mu\phi\,\partial^\mu\phi
  -\frac12m^2\phi^2,
  \qquad
  \partial^\mu=\eta^{\mu\nu}\partial_\nu.
\]
The Euler--Lagrange equation is the Klein--Gordon equation
\[
  (\partial^\mu\partial_\mu-m^2)\phi=0.
\]

Use the Fourier convention
\[
  \phi(x)
  =
  \int\frac{\dd^Dk}{(2\pi)^D}\,
  \widetilde\phi(k)\,\ee^{\ii k\cdot x},
  \qquad
  k\cdot x=\eta_{\mu\nu}k^\mu x^\nu
  =
  -k^0x^0+\vec k\cdot\vec x.
\]
The equation of motion implies
\[
  (k^2+m^2)\widetilde\phi(k)=0,
  \qquad
  k^2=\eta_{\mu\nu}k^\mu k^\nu.
\]
Thus \(\widetilde\phi(k)\) is supported on the mass shell
\[
  k^2+m^2=0,
  \qquad
  k^0=\pm\omega_{\vec k},
  \qquad
  \omega_{\vec k}=\sqrt{\vec k^{\,2}+m^2}.
\]

\begin{center}
\begin{tikzpicture}[>=Stealth,scale=0.95]
  \draw[->] (-3.15,0) -- (3.35,0) node[right] {$|\vec k|$};
  \draw[->] (0,-2.55) -- (0,2.75) node[above] {$k^0$};
  \draw[blue!70!black,line width=1pt,domain=-2.65:2.65,samples=100]
    plot (\x,{sqrt(0.65+\x*\x)});
  \draw[blue!70!black,line width=1pt,domain=-2.65:2.65,samples=100]
    plot (\x,{-sqrt(0.65+\x*\x)});
  \node[blue!70!black,right] at (2.05,2.25) {$k^0=+\omega_{\vec k}$};
  \node[blue!70!black,right] at (2.05,-2.25) {$k^0=-\omega_{\vec k}$};
  \node at (-1.25,0.35) {$k^2+m^2=0$};
\end{tikzpicture}
\end{center}

A real solution can be written as
\[
  \phi(t,\vec x)
  =
  \int\frac{\dd^{D-1}\vec k}{(2\pi)^{D-1}}
  \left[
    A(\vec k)\,\ee^{\ii\vec k\cdot\vec x-\ii\omega_{\vec k}t}
    +
    A(\vec k)^*\,\ee^{-\ii\vec k\cdot\vec x+\ii\omega_{\vec k}t}
  \right],
\]
with a complex coefficient function \(A(\vec k)\) subject to whatever
regularity conditions are imposed on the classical solution. The Cauchy data
\[
  \phi(0,\vec x),\qquad \partial_t\phi(0,\vec x)
\]
determine \(A(\vec k)\) and \(A(\vec k)^*\), hence determine the free
solution for all \(t\). Explicitly, write
\[
  \phi(0,\vec x)
  =
  \int\frac{\dd^{D-1}\vec k}{(2\pi)^{D-1}}\,
  \widetilde\phi_0(\vec k)\ee^{\ii\vec k\cdot\vec x},
  \qquad
  \partial_t\phi(0,\vec x)
  =
  \int\frac{\dd^{D-1}\vec k}{(2\pi)^{D-1}}\,
  \widetilde\pi_0(\vec k)\ee^{\ii\vec k\cdot\vec x}.
\]
Then
\[
  \widetilde\phi_0(\vec k)
  =
  A(\vec k)+A(-\vec k)^*,
  \qquad
  \widetilde\pi_0(\vec k)
  =
  -\ii\omega_{\vec k}\bigl(A(\vec k)-A(-\vec k)^*\bigr),
\]
so that
\[
  A(\vec k)
  =
  \frac12\left(
    \widetilde\phi_0(\vec k)
    +
    \frac{\ii}{\omega_{\vec k}}\widetilde\pi_0(\vec k)
  \right).
\]

There are two constructions to keep distinct. Canonical quantization starts
from the equal-time phase space \((\phi,\Pi)\) and a Hamiltonian. The
Lagrangian path-integral construction is defined by first choosing a regulator
that replaces the field by finitely many coordinates or modes:
\[
  \int_{\mathcal C_R}\dd\mu_R(\phi_R)\,
  \exp\left(\frac{\ii}{\hbar}S_R[\phi_R]\right).
\]
Here \(R\) labels the chosen cutoff, \(\mathcal C_R\) is the regulated
configuration space, \(\dd\mu_R\) is the corresponding finite-dimensional or
lattice measure, and \(S_R\) is the regulated action.  The continuum symbol
\(\int[D\phi]\exp(\ii S[\phi]/\hbar)\) denotes a specified limit or asymptotic
expansion of these regulated objects. This chapter develops the canonical
construction; the scalar field path integral is introduced later.

\section{Canonical Hamiltonian Formalism}

Let
\[
  L(t)=\int\dd^{D-1}\vec x\,\mathcal L.
\]
The canonical momentum density conjugate to \(\phi(t,\vec x)\) is
\[
  \Pi(t,\vec x)
  =
  \frac{\partial\mathcal L}
       {\partial(\partial_t\phi(t,\vec x))}.
\]
For the free scalar,
\[
  \Pi(t,\vec x)=\partial_t\phi(t,\vec x).
\]
The Hamiltonian density is
\[
  \mathcal H
  =
  \Pi\,\partial_t\phi-\mathcal L,
\]
and therefore
\[
  \mathcal H
  =
  \frac12\Pi^2
  +
  \frac12\sum_{i=1}^{D-1}(\partial_i\phi)^2
  +
  \frac12m^2\phi^2
\]
for the free scalar field. The Hamiltonian is
\[
  H=\int\dd^{D-1}\vec x\,\mathcal H.
\]

For functionals \(F[\phi,\Pi]\) and \(G[\phi,\Pi]\), use the equal-time
Poisson bracket
\[
  \{F,G\}
  =
  \int\dd^{D-1}\vec x
  \left[
    \frac{\delta F}{\delta\phi(\vec x)}
    \frac{\delta G}{\delta\Pi(\vec x)}
    -
    \frac{\delta F}{\delta\Pi(\vec x)}
    \frac{\delta G}{\delta\phi(\vec x)}
  \right].
\]
Then
\[
  \{\phi(\vec x),\Pi(\vec y)\}
  =
  \delta^{(D-1)}(\vec x-\vec y),
\]
with the \(\phi\)-\(\phi\) and \(\Pi\)-\(\Pi\) brackets equal to zero. Hamilton's
equation is
\[
  \partial_tF=\{F,H\}.
\]
Applied to \(\phi\) and \(\Pi\), it reproduces the Klein--Gordon equation.

\section{Equal-Time Quantization}

Canonical quantization promotes the classical fields to operator-valued
distributions
\[
  \widehat\phi(t,\vec x),\qquad \widehat\Pi(t,\vec x)
\]
on a common dense domain. Their equal-time commutation relations are
\[
  [\widehat\phi(t,\vec x),\widehat\Pi(t,\vec y)]
  =
  \ii\hbar\,\delta^{(D-1)}(\vec x-\vec y),
\]
and
\[
  [\widehat\phi(t,\vec x),\widehat\phi(t,\vec y)]
  =
  0
  =
  [\widehat\Pi(t,\vec x),\widehat\Pi(t,\vec y)].
\]
All pointwise products are distributional notation. Smeared fields, for
test functions \(f,g\in C_c^\infty(\mathbb R^{D-1})\), are
\[
  \widehat\phi_t(f)=\int\dd^{D-1}\vec x\,f(\vec x)\widehat\phi(t,\vec x),
  \qquad
  \widehat\Pi_t(g)=\int\dd^{D-1}\vec x\,g(\vec x)\widehat\Pi(t,\vec x).
\]
The smeared commutator is
\[
  [\widehat\phi_t(f),\widehat\Pi_t(g)]
  =
  \ii\hbar
  \int\dd^{D-1}\vec x\,f(\vec x)g(\vec x).
\]

For the rest of this chapter set \(\hbar=1\). With \(\hbar\) retained, the
oscillator algebra below reads
\[
  [b_{\vec k},b^\dagger_{\vec q}]
  =
  \hbar(2\pi)^{D-1}\delta^{(D-1)}(\vec k-\vec q).
\]
Equivalently, the \(\hbar=1\) variables are obtained from the \(\hbar\)-full
variables by \(b_{\vec k}^{(\hbar=1)}=\hbar^{-1/2}b_{\vec k}^{(\hbar)}\).

\section{Oscillator Variables and the Hamiltonian}

Let \(\nu_{\vec k}>0\) be a positive even function of \(\vec k\). Define
operator-valued distributions \(b_{\vec k}\) and \(b^\dagger_{\vec k}\) by the
equal-time expansion
\[
  \widehat\phi(0,\vec x)
  =
  \int\frac{\dd^{D-1}\vec k}{(2\pi)^{D-1}}\,
  \frac{1}{\sqrt{2\nu_{\vec k}}}
  \left(
    b_{\vec k}\ee^{\ii\vec k\cdot\vec x}
    +
    b^\dagger_{\vec k}\ee^{-\ii\vec k\cdot\vec x}
  \right),
\]
\[
  \widehat\Pi(0,\vec x)
  =
  \int\frac{\dd^{D-1}\vec k}{(2\pi)^{D-1}}\,
  (-\ii)\sqrt{\frac{\nu_{\vec k}}{2}}
  \left(
    b_{\vec k}\ee^{\ii\vec k\cdot\vec x}
    -
    b^\dagger_{\vec k}\ee^{-\ii\vec k\cdot\vec x}
  \right).
\]
The equal-time commutation relations are equivalent to
\[
  [b_{\vec k},b_{\vec q}]=0
  =
  [b^\dagger_{\vec k},b^\dagger_{\vec q}],
  \qquad
  [b_{\vec k},b^\dagger_{\vec q}]
  =
  (2\pi)^{D-1}\delta^{(D-1)}(\vec k-\vec q).
\]
This algebra holds for every positive choice of \(\nu_{\vec k}\). The
notation \(b_{\vec k}\), \(b^\dagger_{\vec k}\) should therefore not yet be
read as a particle interpretation. If \(\lambda_{\vec k}>0\) is another even
function, define \(c_{\vec k}\), \(c^\dagger_{\vec k}\) by
\[
  c_{\vec k}+c^\dagger_{-\vec k}
  =
  \lambda_{\vec k}\bigl(b_{\vec k}+b^\dagger_{-\vec k}\bigr),
  \qquad
  c_{\vec k}-c^\dagger_{-\vec k}
  =
  \lambda_{\vec k}^{-1}\bigl(b_{\vec k}-b^\dagger_{-\vec k}\bigr).
\]
Equivalently,
\[
  c_{\vec k}
  =
  \frac{\lambda_{\vec k}+\lambda_{\vec k}^{-1}}{2}\,b_{\vec k}
  +
  \frac{\lambda_{\vec k}-\lambda_{\vec k}^{-1}}{2}\,
  b^\dagger_{-\vec k}.
\]
The \(c\)'s obey the same canonical oscillator algebra. If a vector is
annihilated by every \(b_{\vec k}\), it is not annihilated by every
\(c_{\vec k}\) unless \(\lambda_{\vec k}=1\). The Hamiltonian is the datum
that selects the oscillator normalization.

Let
\[
  \omega_{\vec k}=\sqrt{\vec k^{\,2}+m^2}.
\]
Substituting the mode expansion into the free Hamiltonian gives
\[
\begin{aligned}
  \widehat H_0
  =
  \int\frac{\dd^{D-1}\vec k}{(2\pi)^{D-1}}
  \Bigg[
  &\frac{\nu_{\vec k}^2+\omega_{\vec k}^2}{2\nu_{\vec k}}
  \left(
    b^\dagger_{\vec k}b_{\vec k}
    +\frac12(2\pi)^{D-1}\delta^{(D-1)}(0)
  \right)\\
  &+
  \frac{\omega_{\vec k}^2-\nu_{\vec k}^2}{4\nu_{\vec k}}
  \left(
    b_{\vec k}b_{-\vec k}
    +
    b^\dagger_{\vec k}b^\dagger_{-\vec k}
  \right)
  \Bigg].
\end{aligned}
\]
The off-diagonal terms vanish precisely for
\[
  \nu_{\vec k}=\omega_{\vec k}.
\]
With this choice,
\[
  \widehat H_0
  =
  \int\frac{\dd^{D-1}\vec k}{(2\pi)^{D-1}}\,
  \omega_{\vec k}
  \left(
    b^\dagger_{\vec k}b_{\vec k}
    +\frac12(2\pi)^{D-1}\delta^{(D-1)}(0)
  \right).
\]
At finite spatial volume and finite momentum cutoff, the second term is a
finite scalar multiple of the identity. A constant multiple of the identity
does not change Heisenberg evolution, but the quantum Hamiltonian must be a
well-defined self-adjoint operator after the regulator is removed. We
therefore choose the additive energy normalization so that the invariant
vacuum has zero energy. With this normalization the free Hamiltonian is
\[
  \widehat H_0^{\mathrm{ren}}
  =
  \int\frac{\dd^{D-1}\vec k}{(2\pi)^{D-1}}\,
  \omega_{\vec k} b^\dagger_{\vec k}b_{\vec k}.
\]
The free vacuum \(\vac\) is defined by
\[
  b_{\vec k}\vac=0
  \qquad\text{for all }\vec k,
\]
and satisfies \(\widehat H_0^{\mathrm{ren}}\vac=0\).

Time evolution generated by \(\widehat H_0^{\mathrm{ren}}\) gives
\[
  \widehat\phi(x)
  =
  \int\frac{\dd^{D-1}\vec k}{(2\pi)^{D-1}}\,
  \frac{1}{\sqrt{2\omega_{\vec k}}}
  \left(
    b_{\vec k}\ee^{\ii\vec k\cdot\vec x-\ii\omega_{\vec k}t}
    +
    b^\dagger_{\vec k}\ee^{-\ii\vec k\cdot\vec x+\ii\omega_{\vec k}t}
  \right).
\]

The invariant mass-shell normalization used earlier is obtained by defining
\[
  a(\vec k)=\sqrt{2\omega_{\vec k}}\,b_{\vec k},
  \qquad
  a^\dagger(\vec k)=\sqrt{2\omega_{\vec k}}\,b^\dagger_{\vec k}.
\]
Then
\[
  [a(\vec k),a^\dagger(\vec q)]
  =
  (2\pi)^{D-1}2\omega_{\vec k}\,
  \delta^{(D-1)}(\vec k-\vec q),
\]
and
\[
  \widehat\phi(x)
  =
  \int
  \frac{\dd^{D-1}\vec k}{(2\pi)^{D-1}2\omega_{\vec k}}\,
  \left(
    a(\vec k)\ee^{\ii k\cdot x}
    +
    a^\dagger(\vec k)\ee^{-\ii k\cdot x}
  \right),
  \qquad
  k^0=\omega_{\vec k}.
\]
This is the convention used for the free scalar field in the earlier
Fock-space construction.

\section{Covariance and Microcausality}

Let \(\mathcal D_0\) be the finite-particle domain generated by applying
finitely many smeared creation operators to \(\vac\). On \(\mathcal D_0\),
\(\widehat\phi(x)\) is an operator-valued distribution. For
\(f\in C_c^\infty(\mathbb M^D)\),
\[
  \widehat\phi(f)=\int\dd^D x\,f(x)\widehat\phi(x)
\]
is the corresponding smeared operator on \(\mathcal D_0\).

The free Fock representation carries the unitary Poincare representation
induced from the one-particle mass shell. With the above normalization, the
field transforms as a scalar:
\[
  U(a,\Lambda)^{-1}\widehat\phi(x)U(a,\Lambda)
  =
  \widehat\phi(\Lambda x+a).
\]

The commutator is
\[
  [\widehat\phi(x),\widehat\phi(y)]
  =
  \ii\,\Delta(x-y),
\]
where the Pauli--Jordan distribution is
\[
  \Delta(x)
  =
  \frac{1}{\ii}
  \int
  \frac{\dd^{D-1}\vec k}{(2\pi)^{D-1}2\omega_{\vec k}}
  \left(
    \ee^{\ii k\cdot x}
    -
    \ee^{-\ii k\cdot x}
  \right),
  \qquad k^0=\omega_{\vec k}.
\]
This distribution is Lorentz invariant. If \(x\) is spacelike, a Lorentz
transformation sends it to \((0,\vec r)\). At equal time,
\[
  \Delta(0,\vec r)
  =
  \frac{1}{\ii}
  \int
  \frac{\dd^{D-1}\vec k}{(2\pi)^{D-1}2\omega_{\vec k}}
  \left(
    \ee^{\ii\vec k\cdot\vec r}
    -
    \ee^{-\ii\vec k\cdot\vec r}
  \right)
  =
  0,
\]
because the two terms are exchanged by \(\vec k\mapsto-\vec k\). Therefore
\[
  [\widehat\phi(x),\widehat\phi(y)]=0
  \qquad
  \text{when }(x-y)^2>0.
\]
Equivalently,
\[
  [\widehat\phi(f),\widehat\phi(g)]=0
\]
when \(\operatorname{supp}f\) and \(\operatorname{supp}g\) are spacelike
separated. This verifies microcausality for the canonically quantized free
scalar field.

The next step in the construction is to identify the local currents and
charges attached to spacetime symmetries. That is the role of Noether's
theorem in the following chapter.
